%% resumen.tex

% From mitthesis package
% Version: 1.02, 2024/06/19
% Documentation: https://ctan.org/pkg/mitthesis
% Author: David Sanchez-Jacome

\chapter*{Resumen}
\pdfbookmark[0]{Resumen}{resumen}

La demanda continua de mayor rendimiento, eficiencia energética y escalabilidad en los sistemas de 
procesamiento de la información ha impulsado las tecnologías integradas más allá de los límites de 
las arquitecturas puramente electrónicas. Aunque los circuitos integrados electrónicos han 
permitido décadas de avance tecnológico, las restricciones físicas fundamentales relacionadas con 
el ancho de banda, la disipación de potencia y la densidad de interconexión se han convertido en 
cuellos de botella cada vez más críticos. La fotónica integrada programable ha surgido como un 
paradigma convincente para superar estas limitaciones, ofreciendo reconfigurabilidad controlada por 
software sobre una plataforma de hardware común. Sin embargo, la transición de estos procesadores 
versátiles desde los prototipos fundamentales de laboratorio hacia sistemas fabricables a gran 
escala está fuertemente limitada por un "muro de la escalabilidad", un cuello de botella 
multifactorial que abarca las pérdidas ópticas acumulativas, las limitaciones de área física 
(footprint) y las interfaces de control eléctrico de ultra alta densidad.

Para franquear sistemáticamente este muro de la escalabilidad, esta tesis doctoral establece una 
base tecnológica fiable para los circuitos integrados fotónicos programables a gran escala, 
abordando las limitaciones clave en cuatro dimensiones jerárquicas: dispositivo, circuito, 
arquitectura y aplicación a nivel de sistema. A nivel fundamental de dispositivo, se hace uso de la 
optimización de forma por método adjunto (adjoint shape optimization) para diseñar componentes 
fundamentales en silicio sobre aislante (SOI) altamente tolerantes a fallos y capaces de minimizar 
las penalizaciones acumulativas de enrutamiento. Cabe destacar que los interferómetros multimodo 
optimizados lograron una relación de división de potencia casi ideal de 50/50 con picos de exceso 
de pérdidas de tan solo 0,077 dB en toda la banda C extendida. Al integrarse con robustas 
estructuras de adaptación modal (inverted tapers) diseñadas para el acoplamiento lateral (edge 
coupling), estos componentes establecen una columna vertebral física de muy bajas pérdidas capaz de 
soportar topologías de alta densidad.

Basándose en esta capa física optimizada, el salto arquitectónico hacia circuitos a gran escala se 
alcanza mediante la validación tecnológica del encapsulado flip-chip bidimensional. Esto permite 
dejar atrás las limitaciones inherentes al tradicional conexionado eléctrico perimetral 
(wire-bonding). Este cambio de paradigma en el encapsulado ha permitido la integración exitosa y el 
control eléctrico simultáneo de un transceptor inteligente de muy alta densidad compuesto por 1020 
celdas unitarias programables y 332 fotodetectores distribuidos. Para trascender los límites 
funcionales de las mallas de propósito general, esta investigación introduce una topología de 
procesamiento híbrida que integra bloques especializados de alto rendimiento —específicamente, 
redes de difracción en guía de onda (AWG) de 35 canales y filtros de múltiples etapas— dentro de un 
núcleo hexagonal genérico. Además, para resolver los complejos retos que plantea la simulación de 
estas mallas de recirculación, se formula y verifica un novedoso marco de modelado matemático 
inductivo basado en grafos de flujo de señal recursivos.

Finalmente, la viabilidad práctica de esta arquitectura se valida experimentalmente bajo rigurosas 
condiciones de funcionamiento en escenarios de red realistas. La caracterización dinámica del 
subsistema integrado de diversidad de polarización demuestra la capacidad de gestionar de forma 
robusta flujos de datos de alta velocidad a 400G, revelando umbrales estrictos de sensibilidad del 
DSP en el receptor frente al retardo de grupo diferencial (DGD). Ampliando el alcance de la 
evaluación para hacer frente al gran crecimiento del tráfico en los centros de datos, la tesis 
demuestra con éxito el enrutamiento óptico punto a multipunto utilizando un conmutador de circuitos 
ópticos de 32x32 definido por software. De este modo, se logra un multicasting LAN-WDM a 400G sin 
errores con un grado de ramificación (radix) de hasta 1x4. En conjunto, las contribuciones de esta 
investigación validan la capacidad transformadora de la fotónica definida por software, impulsando 
el avance de esta tecnología hacia plataformas robustas y fabricables, capaces de satisfacer las 
demandas dinámicas de ancho de banda y baja latencia de las infraestructuras de telecomunicaciones 
de próxima generación y los centros de datos impulsados por IA.


\chapter*{Resum}
\pdfbookmark[0]{Resum}{resum}

La demanda contínua de major rendiment, eficiència energètica i escalabilitat en els sistemes de 
processament de la informació ha impulsat les tecnologies integrades més enllà dels límits de les 
arquitectures purament electròniques. Encara que els circuits integrats electrònics han permès 
dècades d'avanç tecnològic, les restriccions físiques fonamentals relacionades amb l'ample de 
banda, la dissipació de potència i la densitat d'interconnexió s'han convertit en colls de botella 
cada vegada més crítics. La fotònica integrada programable ha sorgit com un paradigma convincent 
per a superar aquestes limitacions, oferint reconfigurabilitat controlada per programari (software) 
sobre una plataforma de maquinari (hardware) comuna. No obstant això, la transició d'aquests 
processadors versàtils des dels prototips fonamentals de laboratori cap a sistemes fabricables a 
gran escala està fortament limitada per un "mur de l'escalabilitat", un coll de botella 
multifactorial que comprèn les pèrdues òptiques acumulatives, les limitacions d'àrea física 
(footprint) i les interfícies de control elèctric d'ultra alta densitat.

Per a franquejar sistemàticament aquest mur de l'escalabilitat, aquesta tesi doctoral estableix una 
base tecnològica fiable per als circuits integrats fotònics programables a gran escala, abordant 
les limitacions clau en quatre dimensions jeràrquiques: dispositiu, circuit, arquitectura i 
aplicació a nivell de sistema. A nivell fonamental de dispositiu, es fa ús de l'optimització de 
forma per mètode adjunt (adjoint shape optimization) per a dissenyar components fonamentals en 
silici sobre aïllant (SOI) altament tolerants a fallades i capaços de minimitzar les penalitzacions 
acumulatives d'encaminament. Cal destacar que els interferòmetres multimode optimitzats van 
aconseguir una relació de divisió de potència quasi ideal de 50/50 amb pics d'excés de pèrdues de 
tan sols 0,077 dB en tota la banda C estesa. En integrar-se amb robustes estructures d'adaptació 
modal (inverted tapers) dissenyades per a l'acoblament lateral (edge coupling), aquests components 
estableixen una columna vertebral física de molt baixes pèrdues capaç de suportar topologies d'alta 
densitat.

Basant-se en aquesta capa física optimitzada, el salt arquitectònic cap a circuits a gran escala 
s'assoleix mitjançant la validació tecnològica de l'encapsulat flip-chip bidimensional. Això permet 
deixar arrere les limitacions inherents al tradicional conexionat elèctric perimetral 
(wire-bonding). Aquest canvi de paradigma en l'encapsulat ha permès la integració reeixida i el 
control elèctric simultani d'un transceptor intel·ligent de molt alta densitat compost per 1020 
cel·les unitàries programables i 332 fotodetectors distribuïts. Per a transcendir els límits 
funcionals de les malles de propòsit general, aquesta investigació introdueix una topologia de 
processament híbrida que integra blocs especialitzats d'alt rendiment —específicament, xarxes de 
difracció en guia d'ona (AWG) de 35 canals i filtres de múltiples etapes— dins d'un nucli hexagonal 
genèric. A més, per a resoldre els complexos reptes que planteja la simulació d'aquestes malles de 
recirculació, es formula i verifica un innovador marc de modelatge matemàtic inductiu basat en 
grafs de flux de senyal recursius.

Finalment, la viabilitat pràctica d'aquesta arquitectura es valida experimentalment sota rigoroses 
condicions de funcionament en escenaris de xarxa realistes. La caracterització dinàmica del 
subsistema integrat de diversitat de polarització demostra la capacitat de gestionar de forma 
robusta fluxos de dades d'alta velocitat a 400G, revelant llindars estrictes de sensibilitat del 
DSP en el receptor enfront del retard de grup diferencial (DGD). Ampliant l'abast de l'avaluació 
per a fer front al gran creixement del trànsit en els centres de dades, la tesi demostra amb èxit 
l'encaminament òptic punt a multipunt utilitzant un commutador de circuits òptics de 32x32 definit 
per programari. D'aquesta manera, s'aconsegueix un multicasting LAN-WDM a 400G sense errors amb un 
grau de ramificació (radix) de fins a 1x4. En conjunt, les contribucions d'aquesta investigació 
validen la capacitat transformadora de la fotònica definida per programari, impulsant l'avanç 
d'aquesta tecnologia cap a plataformes robustes i fabricables, capaces de satisfer les demandes 
dinàmiques d'ample de banda i baixa latència de les infraestructures de telecomunicacions de 
pròxima generació i els centres de dades impulsats per IA.
